
\documentclass{article}
\usepackage{colortbl}
\usepackage{makecell}
\usepackage{multirow}
\usepackage{supertabular}

\begin{document}

\newcounter{utterance}

\centering \large Interaction Transcript for game `eqbench', experiment `eq\_bench', episode 163 with reflexion{-}Qwen3.5{-}27B.
\vspace{24pt}

{ \footnotesize  \setcounter{utterance}{1}
\setlength{\tabcolsep}{0pt}
\begin{supertabular}{c@{$\;$}|p{.15\linewidth}@{}p{.15\linewidth}p{.15\linewidth}p{.15\linewidth}p{.15\linewidth}p{.15\linewidth}}
    \# & \multicolumn{2}{c}{Player} && \multicolumn{2}{c}{Game Master} \\
    \hline

    \theutterance \stepcounter{utterance}  
    & & & \multicolumn{4}{p{0.6\linewidth}}{
        \cellcolor[rgb]{0.9,0.9,0.9}{
            \makecell[{{p{\linewidth}}}]{
                \texttt{\tiny{[P1$\langle$GM]}}
                \texttt{You have attempted this same task before and did not succeed. These are the plans you wrote after those attempts.} \\
\texttt{Trial \#1: New plan: My previous attempt failed because I assigned a high score to "Intrigued" (8) while ignoring the emotional weight of Jonny's final comment about loneliness. The dialogue shifts from playful banter to a moment of vulnerability where Jonny explicitly states, "Must be lonely," and Eva's response ("Something like that") implies a shared, heavy realization rather than mere curiosity. I should have prioritized "Misunderstood" or "Pity" (for herself or others) over "Intrigued." Next time, I will analyze the final line of the dialogue as the primary emotional anchor. If the dialogue ends with a statement about isolation or sadness, I will assign the highest score to the emotion reflecting that specific sentiment (e.g., Misunderstood or Pity) and lower the score for lighter emotions like Intrigued, ensuring the scores reflect the narrative arc's conclusion rather than the middle of the conversation.} \\
\\ 
\texttt{Trial \#2: New plan: My previous attempt failed because I overestimated "Misunderstood" and "Pity" while underestimating the positive engagement in the dialogue. Jonny's comment about loneliness is met with Eva's laughter and a calm admission of her nature, suggesting she is comfortable with her difference rather than feeling isolated or pitied. The interaction is characterized by a unique connection where Jonny validates her, making "Intrigued" the most accurate primary emotion as she engages with someone who finally "sees" her. Next time, I will re{-}evaluate the tone of the final exchange: if the character responds to a heavy topic with humor or calm acceptance, I will lower scores for negative emotions like "Misunderstood" and "Pity" and significantly raise the score for "Intrigued" to reflect the successful connection and curiosity in the relationship. I will ensure the highest score aligns with the character's active, positive engagement rather than passive suffering.} \\
\\ 
\\ 
\\ 
\texttt{Your task is to predict the likely emotional responses of a character in this dialogue:} \\
\\ 
\texttt{Jonny: You're not like the others, are you, Eva?} \\
\texttt{Eva: And what makes you say that?} \\
\texttt{Jonny: Well, for starters, you can see me. And you don't scream and run away.} \\
\texttt{Eva: (laughs) I suppose that does make me a bit different.} \\
\texttt{Jonny: More than a bit, I'd say. } \\
\texttt{Eva: I've always been... sensitive to things others can't see. } \\
\texttt{Jonny: Like a sixth sense?} \\
\texttt{Eva: Something like that. } \\
\texttt{Jonny: Must be lonely, seeing things no one else can.} \\
\texttt{{[}End dialogue{]}} \\
\\ 
\texttt{At the end of this dialogue, Eva would feel...} \\
\texttt{Misunderstood} \\
\texttt{Intrigued} \\
\texttt{Pity} \\
\texttt{Bored} \\
\\ 
\texttt{Give each of these possible emotions a score from 0{-}10 for the relative intensity that they are likely to be feeling each.} \\
\\ 
\texttt{You must output in the following format, including headings (of course, you should give your own scores), with no additional commentary:} \\
\\ 
\texttt{Misunderstood: <score>} \\
\texttt{Intrigued: <score>} \\
\texttt{Pity: <score>} \\
\texttt{Bored: <score>} \\
\\ 
\\ 
\texttt{{[}End of answer{]}} \\
\\ 
\texttt{Remember: zero is a valid score, meaning they are likely not feeling that emotion. You must score at least one emotion > 0.} \\
\\ 
\texttt{Your answer:} \\
            }
        }
    }
     \\ \\

    \theutterance \stepcounter{utterance}  
    & \multicolumn{4}{p{0.6\linewidth}}{
        \cellcolor[rgb]{0.561,0.737,0.561}{
            \makecell[{{p{\linewidth}}}]{
                \texttt{\tiny{[P1$\rangle$GM]}}
                \texttt{Misunderstood: 2} \\
\texttt{Intrigued: 8} \\
\texttt{Pity: 1} \\
\texttt{Bored: 0} \\
            }
        }
    }
    & & \\ \\

    \theutterance \stepcounter{utterance}  
    & & & \multicolumn{2}{p{0.3\linewidth}}{
        \cellcolor[rgb]{0.9,0.9,0.9}{
            \makecell[{{p{\linewidth}}}]{
                \texttt{\tiny{[GM$|$GM]}}
                \texttt{Misunderstood: 2} \\
\texttt{Intrigued: 8} \\
\texttt{Pity: 1} \\
\texttt{Bored: 0} \\
            }
        }
    }
    & & \\ \\

    \theutterance \stepcounter{utterance}  
    & & & \multicolumn{2}{p{0.3\linewidth}}{
        \cellcolor[rgb]{0.9,0.9,0.9}{
            \makecell[{{p{\linewidth}}}]{
                \texttt{\tiny{[GM$|$GM]}}
                \texttt{Misunderstood: 2} \\
\texttt{Intrigued: 8} \\
\texttt{Pity: 4} \\
\texttt{Bored: 0} \\
            }
        }
    }
    & & \\ \\

    \theutterance \stepcounter{utterance}  
    & & & \multicolumn{2}{p{0.3\linewidth}}{
        \cellcolor[rgb]{0.9,0.9,0.9}{
            \makecell[{{p{\linewidth}}}]{
                \texttt{\tiny{[GM$|$GM]}}
                \texttt{game\_result = LOSE} \\
            }
        }
    }
    & & \\ \\

    \theutterance \stepcounter{utterance}  
    & & \multicolumn{4}{p{0.6\linewidth}}{
        \cellcolor[rgb]{0.95,0.95,0.95}{
            \makecell[{{p{\linewidth}}}]{
                \texttt{\tiny{[P1$\leftrightarrow$P]}}
                \texttt{Post{-}episode Reflexion (not sent to GM):} \\
\texttt{New plan: My previous attempts failed because I oscillated between interpreting the dialogue as purely melancholic (Trial 1) or purely positive (Trial 2), missing the nuanced blend of emotions. The key is to recognize that Eva's laughter and admission ("Something like that") in response to the "lonely" comment indicate a complex state: she is **Intrigued** by Jonny's unique perception of her, but the topic of loneliness inherently triggers a sense of being **Misunderstood** by the general world, even if Jonny understands her. I should not treat these as mutually exclusive. Next time, I will assign the highest score to "Intrigued" (reflecting the active connection) but also assign a significant, non{-}zero score to "Misunderstood" (reflecting the underlying theme of isolation she acknowledges), while keeping "Pity" and "Bored" low. I will ensure the scores reflect this duality: high engagement with the specific person, but an acknowledgment of the general condition of being different.} \\
            }
        }
    }
    & \\ \\

\end{supertabular}
}

\end{document}
